\chapter{使用示例}

本章给出一套适合当前模板的图表写作规范。推荐原则是：
普通插图使用 \pkg{graphicx}，子图统一使用 \pkg{subcaption}，
结构示意图使用 \pkg{tikz}，数据图优先使用 \pkg{pgfplots}；
简单表格使用 \env{tabular}+\pkg{booktabs}，
自动列宽表格使用 \env{tabularx}，
跨页长表使用 \env{longtable}，
复杂样式表格优先使用 \pkg{tabularray}。

\section{插图}

图形部分建议采用下面这套固定工作流：
\begin{enumerate}
\item 普通插图优先使用 \verb|PDF|、\verb|SVG| 转换后的矢量图；
\item 子图统一使用 \pkg{subcaption}，不要再混用旧的 \pkg{subfig}；
\item 结构示意图、流程图、框图使用 \pkg{tikz}；
\item 实验曲线、柱状图、三维曲面图等数据图优先使用 \pkg{pgfplots}；
\item 只有确实需要整页横置时，才使用 \pkg{rotating} 提供的 \env{sidewaysfigure}。
\end{enumerate}

\begin{figure}[tbp]
  \centering
  \newcounter{density}
  \setcounter{density}{20}
  \input{fig/tikz_rot}
  \caption{使用 \pkg{tikz} 绘制的结构示意图}
  \label{fig:tikzrot}
\end{figure}

\begin{figure}[tbp]
  \centering
  \includegraphics[width=4cm]{nuaa-logo.pdf}
  \caption{外部矢量图示例}
  \label{fig:logo}
\end{figure}

如果需要多张图片共用一个总题注，推荐直接使用 \pkg{subcaption}，如\autoref{fig:sub2}。

\begin{figure}[tbp]
  \centering
  \begin{subfigure}[b]{0.48\textwidth}
    \centering
    \includegraphics[width=4cm]{nuaa-logo.pdf}
    \caption{左边的大校徽}
    \label{fig:sub1}
  \end{subfigure}
  \hfill
  \begin{subfigure}[b]{0.40\textwidth}
    \centering
    \includegraphics[width=3cm]{nuaa-logo.pdf}
    \caption{小校徽，题注较长时会自动换行}
    \label{fig:sub2}
  \end{subfigure}
  \caption{使用 \pkg{subcaption} 的子图示例}
  \label{fig:subfigs}
\end{figure}

如果需要插入实验数据图，可以优先考虑 \pkg{pgfplots}。这类图的代码结构更清晰，
也更容易和全文字体、字号统一，效果见\autoref{fig:plots}。

\begin{figure}[tbp]
  \centering
  \begin{subfigure}[b]{0.48\textwidth}
    \centering
    \input{fig/plot_2d}
    \caption{二维折线图}
    \label{fig:func}
  \end{subfigure}
  \hfill
  \begin{subfigure}[b]{0.48\textwidth}
    \centering
    \input{fig/plot_3d}
    \caption{三维曲面图}
    \label{fig:sum}
  \end{subfigure}
  \caption{使用 \pkg{pgfplots} 绘制的数据图}
  \label{fig:plots}
\end{figure}

如果需要让多张图片分两页共用一个题注，可以继续使用 \cs{ContinuedFloat}，
并在续页标题里使用空的可选参数，避免图表清单重复条目。

\begin{figure}[t]
  \centering
  \begin{subfigure}[b]{0.42\textwidth}
    \centering
    \includegraphics[width=4cm]{nuaa-logo.pdf}
    \caption{校徽$\times 1$}
  \end{subfigure}
  \hfill
  \begin{subfigure}[b]{0.42\textwidth}
    \centering
    \includegraphics[width=.4\textwidth]{nuaa-logo.pdf}
    \caption{校徽$\times 2$}
  \end{subfigure}

  \medskip

  \begin{subfigure}[b]{0.42\textwidth}
    \centering
    \includegraphics[width=.4\textwidth]{nuaa-logo.pdf}
    \caption{校徽$\times 3$}
  \end{subfigure}
  \hfill
  \begin{subfigure}[b]{0.42\textwidth}
    \centering
    \includegraphics[width=4cm]{nuaa-logo.pdf}
    \caption{校徽$\times 4$}
  \end{subfigure}
  \caption{多图分页示例}
  \label{fig:subfigss}
\end{figure}

\begin{figure}[t]
  \ContinuedFloat
  \centering
  \begin{subfigure}[b]{0.42\textwidth}
    \centering
    \includegraphics[width=4cm]{nuaa-logo.pdf}
    \caption{校徽$\times 5$}
  \end{subfigure}
  \hfill
  \begin{subfigure}[b]{0.42\textwidth}
    \centering
    \includegraphics[width=4cm]{nuaa-logo.pdf}
    \caption{校徽$\times 6$}
    \label{fig:logo6}
  \end{subfigure}

  \medskip

  \begin{subfigure}[b]{0.42\textwidth}
    \centering
    \includegraphics[width=4cm]{nuaa-logo.pdf}
    \caption{校徽$\times 7$}
  \end{subfigure}
  \hfill
  \begin{subfigure}[b]{0.42\textwidth}
    \centering
    \includegraphics[width=4cm]{nuaa-logo.pdf}
    \caption{校徽$\times 8$}
  \end{subfigure}
  \caption[]{多图分页示例（续）}
\end{figure}

如果需要将大图横置为单独页面，建议使用 \env{sidewaysfigure}，
它比手工旋转整页图片更稳定。

\begin{sidewaysfigure}
  \centering
  \includegraphics[width=.82\textheight]{nuaa-jianqi.pdf}
  \caption{整页横置图片示例}
  \label{fig:fullpage}
\end{sidewaysfigure}

\section{表格}

表格部分建议采用下面这套固定工作流：
\begin{enumerate}
\item 简单表格使用 \env{tabular}+\pkg{booktabs}；
\item 列内容较长时使用 \env{tabularx} 做自动列宽；
\item 跨页长表使用 \env{longtable}；
\item 复杂样式或需要统一控制边框、列格式时，优先使用 \pkg{tabularray}。
\end{enumerate}

如果需要插入一个简易表格，可以只使用 \env{tabular} 环境，如\autoref{tab:city}。
\begin{table}[tbp]
  \centering
  \caption[城市人口]{城市人口数量排名（source: Wikipedia）\label{tab:city}}
  \begin{tabular}{lr}
    \toprule
    城市 & 人口 \\
    \midrule
    Mexico City & 20,116,842 \\
    Shanghai & 19,210,000 \\
    Beijing & 15,796,450 \\
    Istanbul & 14,160,467 \\
    \bottomrule
  \end{tabular}
\end{table}

如果列内容较长，推荐使用 \env{tabularx} 自动分配剩余宽度，而不是继续使用老旧的 \env{tabu}。

\begin{table}[tbp]
  \centering
  \caption{\env{tabularx} 自动列宽示例 \label{tab:tabx}}
  \begin{tabularx}{.92\textwidth}{@{}Y Y@{}}
    \toprule
    组件 & 说明 \\
    \midrule
    \pkg{graphicx} & 用于插入外部图片，是普通插图的基础工具。 \\
    \pkg{subcaption} & 用于组织子图和子表，适合多图并排对比与统一编号。 \\
    \pkg{pgfplots} & 适合绘制实验曲线、柱状图和散点图，便于统一全文风格。 \\
    \pkg{tabularx} & 适合正文里说明文字较多、需要自动列宽的表格。 \\
    \bottomrule
  \end{tabularx}
\end{table}

如果需要更现代的表格样式控制，可以使用 \pkg{tabularray} 的 \env{tblr} 环境。

\begin{table}[tbp]
  \centering
  \caption{\pkg{tabularray} 示例 \label{tab:tblr}}
  \begin{tblr}{
    colspec = {Q[l] Q[l] Q[c]},
    row{1} = {font=\bfseries},
    hlines,
    vlines
  }
    模块 & 说明 & 推荐场景 \\
    \pkg{graphicx} & 外部图片插入 & 矢量图、位图 \\
    \pkg{subcaption} & 子图与子表编排 & 多图并排、对比实验 \\
    \pkg{pgfplots} & 数据图绘制 & 曲线图、柱状图、曲面图 \\
    \pkg{longtable} & 跨页表格 & 大型实验数据表 \\
  \end{tblr}
\end{table}

如果需要对某一列的小数点对齐、带单位或做四舍五入处理，可以配合 \pkg{siunitx} 使用。
\autoref{tab:xmpl:mixed} 改编自该宏包文档中的数值与单位排版示例。

\begin{table}[tbp]
  \centering
  \caption{Tables where numbers have different units}
  \label{tab:xmpl:mixed}
  \begin{tabular}
    {
      @{}
      >{$}l<{$}
      S[table-format = 3.3(1)]
      S[table-format = 3.3(1)]
      @{}
    }
    \toprule
      & {One} & {Two} \\
    \midrule
    a / \unit{\pm}         & 123.4(2) & 567.8(4) \\
    \beta / \unit{\degree} & 90.34(4) & 104.45(5) \\
    \mu / \unit{\per\mm}   & 0.532    & 0.894    \\
    \bottomrule
  \end{tabular}
  \hfil
  \begin{tabular}
    {
      @{}
      S[table-format=1.3]@{\,}
      >{\collectcell\unit}l<{\endcollectcell}
      @{}
    }
    \toprule
    \multicolumn{2}{@{}c}{Heading} \\
    \midrule
    1.234 & \metre \\
    0.835 & \candela \\
    4.23  & \joule\per\mole \\
    \bottomrule
  \end{tabular}
\end{table}

如果表格内容很多，导致无法放在一页内，应使用 \env{longtable} 进行分页。
\autoref{tab:performance} 展示了一个跨页长表的基本写法。

\begin{longtable}[c]{@{}c*{6}{r}@{}}
  \caption[实验数据]{实验数据，这个题注十分的长，注意这在索引中的处理方式，还有 \cs{caption} 后面的双反斜杠}\label{tab:performance}\\
  \toprule
  \multirow{2}{*}{测试程序} & \multicolumn{1}{c}{正常运行} & \multicolumn{1}{c}{同步} & \multicolumn{1}{c}{检查点} & \multicolumn{1}{c}{卷回恢复} &
  \multicolumn{1}{c}{进程迁移} & \multicolumn{1}{c}{检查点} \\
  & \multicolumn{1}{c}{时间 (s)} & \multicolumn{1}{c}{时间 (s)} &
  \multicolumn{1}{c}{时间 (s)} & \multicolumn{1}{c}{时间 (s)} &
  \multicolumn{1}{c}{时间 (s)} & \multicolumn{1}{c}{文件 (KB)} \\
  \midrule
  \endfirsthead

  \multicolumn{7}{c}{\nuaafontcaption 表~\thetable（续）}\\
  \toprule
  \multirow{2}{*}{测试程序} & \multicolumn{1}{c}{正常运行} & \multicolumn{1}{c}{同步} & \multicolumn{1}{c}{检查点} & \multicolumn{1}{c}{卷回恢复} &
  \multicolumn{1}{c}{进程迁移} & \multicolumn{1}{c}{检查点} \\
  & \multicolumn{1}{c}{时间 (s)} & \multicolumn{1}{c}{时间 (s)} &
  \multicolumn{1}{c}{时间 (s)} & \multicolumn{1}{c}{时间 (s)} &
  \multicolumn{1}{c}{时间 (s)} & \multicolumn{1}{c}{文件 (KB)} \\
  \midrule
  \endhead

  \midrule
  \multicolumn{7}{r}{续下页}
  \endfoot

  \bottomrule
  \endlastfoot

  CG.A.2 & 23.05 & 0.002 & 0.116 & 0.035 & 0.589 & 32491 \\
  CG.A.4 & 15.06 & 0.003 & 0.067 & 0.021 & 0.351 & 18211 \\
  CG.A.8 & 13.38 & 0.004 & 0.072 & 0.023 & 0.210 & 9890 \\
  CG.B.2 & 867.45 & 0.002 & 0.864 & 0.232 & 3.256 & 228562 \\
  CG.B.4 & 501.61 & 0.003 & 0.438 & 0.136 & 2.075 & 123862 \\
  CG.B.8 & 384.65 & 0.004 & 0.457 & 0.108 & 1.235 & 63777 \\
  MG.A.2 & 112.27 & 0.002 & 0.846 & 0.237 & 3.930 & 236473 \\
  MG.A.4 & 59.84 & 0.003 & 0.442 & 0.128 & 2.070 & 123875 \\
  MG.A.8 & 31.38 & 0.003 & 0.476 & 0.114 & 1.041 & 60627 \\
  MG.B.2 & 526.28 & 0.002 & 0.821 & 0.238 & 4.176 & 236635 \\
  MG.B.4 & 280.11 & 0.003 & 0.432 & 0.130 & 1.706 & 123793 \\
  MG.B.8 & 148.29 & 0.003 & 0.442 & 0.116 & 0.893 & 60600 \\
  LU.A.2 & 2116.54 & 0.002 & 0.110 & 0.030 & 0.532 & 28754 \\
  LU.A.4 & 1102.50 & 0.002 & 0.069 & 0.017 & 0.255 & 14915 \\
  LU.A.8 & 574.47 & 0.003 & 0.067 & 0.016 & 0.192 & 8655 \\
  LU.B.2 & 9712.87 & 0.002 & 0.357 & 0.104 & 1.734 & 101975 \\
  LU.B.4 & 4757.80 & 0.003 & 0.190 & 0.056 & 0.808 & 53522 \\
  LU.B.8 & 2444.05 & 0.004 & 0.222 & 0.057 & 0.548 & 30134 \\
  CG.B.2 & 867.45 & 0.002 & 0.864 & 0.232 & 3.256 & 228562 \\
  CG.B.4 & 501.61 & 0.003 & 0.438 & 0.136 & 2.075 & 123862 \\
  CG.B.8 & 384.65 & 0.004 & 0.457 & 0.108 & 1.235 & 63777 \\
  MG.A.2 & 112.27 & 0.002 & 0.846 & 0.237 & 3.930 & 236473 \\
  MG.A.4 & 59.84 & 0.003 & 0.442 & 0.128 & 2.070 & 123875 \\
  MG.A.8 & 31.38 & 0.003 & 0.476 & 0.114 & 1.041 & 60627 \\
  MG.B.2 & 526.28 & 0.002 & 0.821 & 0.238 & 4.176 & 236635 \\
  MG.B.4 & 280.11 & 0.003 & 0.432 & 0.130 & 1.706 & 123793 \\
  MG.B.8 & 148.29 & 0.003 & 0.442 & 0.116 & 0.893 & 60600 \\
  LU.A.2 & 2116.54 & 0.002 & 0.110 & 0.030 & 0.532 & 28754 \\
  LU.A.4 & 1102.50 & 0.002 & 0.069 & 0.017 & 0.255 & 14915 \\
  LU.A.8 & 574.47 & 0.003 & 0.067 & 0.016 & 0.192 & 8655 \\
  LU.B.2 & 9712.87 & 0.002 & 0.357 & 0.104 & 1.734 & 101975 \\
  LU.B.4 & 4757.80 & 0.003 & 0.190 & 0.056 & 0.808 & 53522 \\
  LU.B.8 & 2444.05 & 0.004 & 0.222 & 0.057 & 0.548 & 30134 \\
  EP.A.2 & 123.81 & 0.002 & 0.010 & 0.003 & 0.074 & 1834 \\
  EP.A.4 & 61.92 & 0.003 & 0.011 & 0.004 & 0.073 & 1743 \\
  EP.A.8 & 31.06 & 0.004 & 0.017 & 0.005 & 0.073 & 1661 \\
  EP.B.2 & 495.49 & 0.001 & 0.009 & 0.003 & 0.196 & 2011 \\
  EP.B.4 & 247.69 & 0.002 & 0.012 & 0.004 & 0.122 & 1663 \\
  EP.B.8 & 126.74 & 0.003 & 0.017 & 0.005 & 0.083 & 1656 \\
\end{longtable}

\section{数字与国际单位}

本模板预加载 \pkg{siunitx} 来格式化文中的内联数字，该宏包有大量可定制的参数，
请务必阅读其文档，并在文档导言部分设置格式。

\begin{itemize}
  \item 旋转角度为 \ang{90}、\ang{270}
  \item 分辨率 \numproduct{1920x1080} 的像素数量约为 \num{2.07e6}
  \item 电脑显示器的像素间距为 \SI{1.8}{\nm}、\SI{180}{\um} 还是 \SI{18}{\mm}？
  \item 重力加速度 $g=\SI{9.8}{\kg\per\square\second}$、
  $g=\SI[inter-unit-product=\ensuremath{{}\cdot{}}]{9.8}{\kg\per\square\second}$，
  亦或 $g=\SI[per-mode=symbol]{9.8}{\kg\per\square\second}$
\end{itemize}

\section{中英文之间空格}

很遗憾，目前 \LaTeX{} 和 \CTeX{} 虽然能处理普通汉字与英文之间的间隔，
但是汉字与宏之间的空格仍然需要手工调整，请务必按以下的规则撰写原稿：
\begin{itemize}
  \item[\ding{51}] 如\autoref{fig:sub2} 所示：\verb|如\autoref{fig:sub2} 所示|，这个宏返回的是“图 x.xx”，
  所以前面两个汉字之间不能加空格，后面数字与汉字之间必须加空格；
  \item[\ding{51}] 距离为 1.7~个天文单位：\verb|距离为 1.7~个天文单位|，前面可以不加空格（\CTeX 会修正），
  后面必须加 \verb|~| 以防止在 “1.7”与“个”之间换行。此时更推荐写成 \SI{1.7}{au}：\verb|\SI{1.7}{au}|。
\end{itemize}
